% pauli_exclusion_ML.tex
% "Pauli's Heuristic as Statistical Inference: The Exclusion Principle
%  as a Parsimonious Four-Parameter Model"
% Thomas J. Sargent, 2026

\documentclass[12pt]{article}

\usepackage{amsmath,amssymb}
\usepackage{booktabs}
\usepackage{array}
\usepackage{multirow}
\usepackage[margin=1.2in]{geometry}
\usepackage{setspace}
\usepackage{microtype}
\usepackage{hyperref}
\usepackage{natbib}
\usepackage{parskip}

\onehalfspacing

\hypersetup{
  colorlinks  = true,
  linkcolor   = blue,
  citecolor   = blue,
  urlcolor    = blue
}

\title{Pauli's Heuristic as Statistical Inference:\\
       The Exclusion Principle as a Parsimonious\\
       Four-Parameter Model, 1913--1928}
\author{Thomas J.\ Sargent}
\date{April 2026}

\begin{document}
\maketitle

\begin{abstract}
In his 1925 paper on the ``complex structure of spectra,'' Wolfgang Pauli
proposed that each electron in an atom is characterised by exactly four quantum
numbers---the principal quantum number~$n$, the azimuthal (orbital) quantum
number~$l$, the magnetic quantum number~$m_l$, and a newly introduced
fourth, ``two-valued'' quantum number~$m_s$---and that no two electrons in any
atom may share the same 4-tuple $(n,l,m_l,m_s)$.  This note examines Pauli's
reasoning from a modern statistical perspective.  The empirical regularities
that Pauli organised---electron shell capacities of 2, 8, 18, 32; the doublet
structure of alkali spectra; the anomalous Zeeman effect; the X-ray
spectroscopic data systematised by Stoner (1924)---constitute a discrete
dataset of counts and multiplicities.  Pauli's four-parameter state space
predicts those counts \emph{exactly} via the formula: capacity of subshell~$l$
equals $2(2l{+}1)$, implying a total shell capacity of $2n^2$.  We frame this
as a highly constrained statistical model in which the parameters are the
allowed ranges of four discrete quantum numbers and a binary occupation
constraint, and we examine which ``data'' identified each parameter.  Pauli's
procedure parallels Einstein's photoelectric derivation: theoretical insight
restricted the hypothesis class to a parsimonious form, reducing what might
have been an unmanageable combinatorial estimation problem to a counting
exercise whose predictions match every relevant observation.  The episode
illustrates how a one-sentence theoretical postulate---the Exclusion
Principle---functions as an extreme form of inductive bias, compressing
the effective hypothesis space from an open-ended combinatorial family to
a single, fully determined prediction with zero tunable parameters.
\end{abstract}

\tableofcontents

%======================================================================
\section{Introduction}

In December 1900, Max Planck introduced the energy quantum $h\nu$ to fit the
blackbody spectrum.  Five years later, Einstein applied the same quantum to the
photoelectric effect, deriving a two-parameter linear equation
$V_0 = (h/e)\nu - W/e$ whose slope encoded Planck's constant.  Twenty years
after that, in January 1925, Wolfgang Pauli proposed something structurally
similar: a minimal theoretical principle from which a large body of empirical
facts about atomic spectra and the periodic table follows as a counting
exercise.

Pauli's principle---that no two electrons in an atom may occupy the same
quantum state---is presented in textbooks as a law of nature, delivered with
an air of inevitability.  The historical record tells a more interesting story.
Pauli \emph{conjectured} the principle as the simplest model consistent with a
specific collection of quantitative empirical regularities that had accumulated
between 1913 and 1924.  He called his fourth quantum number a
``classically non-describable two-valuedness''~[\textit{eine nicht klassisch
beschreibbare Zweideutigkeit}] because he had no physical interpretation for
it.  It was, in his own words, a formal device introduced ``at any cost'' to
organise the data.  In this respect, Pauli's reasoning is a closer relative of
Planck's 1900 ``act of desperation'' than of the theoretical deductions from
first principles that characterise later quantum mechanics.

This note examines Pauli's 1925 paper through the lens of modern statistical
inference.  The central questions are:
\begin{enumerate}
  \item What were the empirical \emph{data} that Pauli was organising?
  \item What is the four-parameter \emph{model} that his principle defines,
    and how many distinct states does it predict for each shell and subshell?
  \item Which parameter---or which constraint---is identified by which
    observation?
  \item How does the episode compare, as an exercise in inference and
    hypothesis formation, with Einstein's derivation of the photoelectric
    equation?
\end{enumerate}
Sections~\ref{sec:history}--\ref{sec:stoner} sketch the experimental history
from Bohr's 1913 model to Stoner's 1924 consolidation.
Section~\ref{sec:model} defines Pauli's four-parameter state space precisely
and works through the counting algebra.
Section~\ref{sec:data} assembles the empirical ``dataset''---the shell
capacities, spectral multiplicities, and anomalous Zeeman observations---that
the model is required to fit.
Section~\ref{sec:fit} examines the goodness of fit and the identification of
each parameter from distinct empirical observations.
Section~\ref{sec:identification} discusses the identification logic: which
data identify each parameter.
Section~\ref{sec:mle} frames the counting problem as a maximum likelihood
estimation problem over a discrete state space.
Section~\ref{sec:subsequent} traces the subsequent confirmations (Uhlenbeck and
Goudsmit's spin hypothesis; Fermi--Dirac statistics; Dirac's relativistic
derivation; the spin--statistics theorem).
Section~\ref{sec:ml} reinterprets the episode in the language of machine
learning.
Section~\ref{sec:comparison} compares Pauli's episode with Einstein's
photoelectric derivation.
Section~\ref{sec:conclusions} concludes.

%======================================================================
\section{Historical Context, 1913--1924}
\label{sec:history}

\subsection{The Bohr Model and Its Three Quantum Numbers}

Niels Bohr published his model of the hydrogen atom in 1913
\citep{Bohr1913}.  The key postulate was that the electron's orbital angular
momentum is quantised: $L = n\hbar$, $n=1,2,3,\ldots$  This single integer,
the \emph{principal quantum number}~$n$, gives the energy levels
\begin{equation}\label{eq:bohr}
  E_n = -\frac{13.6\,\text{eV}}{n^2},
\end{equation}
correctly predicting the hydrogen emission spectrum (Lyman, Balmer, and
Paschen series) to parts per thousand.

Arnold Sommerfeld extended Bohr's model in 1916~\citep{Sommerfeld1916} by
allowing elliptical orbits and applying the Wilson--Sommerfeld quantisation
condition $\oint p_i\,dq_i = n_i h$ to each coordinate.  This introduced two
additional integers:
\begin{itemize}
  \item The \emph{azimuthal quantum number}~$l$ (Sommerfeld's $k$,
    relabelled later): $l\in\{0,1,\ldots,n-1\}$, determining the orbital
    angular momentum magnitude $|\mathbf{L}| = \sqrt{l(l+1)}\,\hbar$.
  \item The \emph{magnetic quantum number}~$m_l$: $m_l\in\{-l,-l{+}1,\ldots,
    +l\}$, giving the projection $L_z = m_l\hbar$ of orbital angular momentum
    on an external magnetic field axis.  In a field, each $m_l$ value produces
    a slightly different energy (Zeeman splitting).
\end{itemize}

With these three integers $(n,l,m_l)$ the number of distinct states in shell~$n$
is
\begin{equation}\label{eq:three_count}
  \sum_{l=0}^{n-1}(2l+1) = n^2.
\end{equation}
Table~\ref{tab:three_param} lists these counts for the first four shells.

\begin{table}[ht]
\centering
\caption{States per shell under the three-quantum-number
Bohr--Sommerfeld model vs.\ observed electron capacities.
The observed capacity equals $2n^2$---exactly twice the predicted
$n^2$---a discrepancy that motivated Pauli's fourth quantum number.}
\label{tab:three_param}
\medskip
\begin{tabular}{cccc}
\toprule
Shell & $n$ & Predicted states ($n^2$) & Observed capacity ($2n^2$) \\
\midrule
K & 1 &  1 &  2 \\
L & 2 &  4 &  8 \\
M & 3 &  9 & 18 \\
N & 4 & 16 & 32 \\
\bottomrule
\end{tabular}
\end{table}

\noindent The predicted capacities fall short by a factor of exactly~2.
This was known but unexplained before 1925.

\subsection{Failures of the Three-Parameter Model}

Despite its successes for hydrogen, the Bohr--Sommerfeld three-quantum-number
scheme faced three acute empirical failures by 1924.

\paragraph{Shell closures and the periodic table.}
Bohr proposed in 1922 that the chemical periodic table is organised by fills
of successive electron shells, with noble gases corresponding to completed
shells: helium has 2 electrons (shell $n=1$ full), neon has 10 (shells $n=1,2$
full), argon has 18 ($n=1,2,3$ partially full in the $3s,3p$ subshells), krypton
has 36, and xenon has 54.  No dynamical principle within the
three-quantum-number scheme explained why a shell can hold exactly $2n^2$
electrons---rather than $n^2$, or $3n^2$, or any other number---or why it is
energetically ``closed'' once filled.

\paragraph{The anomalous Zeeman effect.}
Pieter Zeeman discovered in 1896--97 \citep{Zeeman1897} that magnetic fields
split spectral lines.  Classical electron theory (Lorentz) predicted the
\emph{normal Zeeman effect}: each line splits into exactly three components
($\Delta m_l = 0,\pm 1$).  In reality, most lines in multi-electron atoms
split into \emph{more} than three components, with fractional-unit spacings.
This \emph{anomalous Zeeman effect} was inexplicable within the
three-quantum-number scheme.  Pauli in 1924 wrote the most famous frustrated
lament in the history of atomic physics, reportedly saying:
\begin{quote}
  ``How can one look happy when he is thinking about the anomalous Zeeman
  effect?'' \citep{Pais1991}
\end{quote}

\paragraph{Half-integer angular momenta.}
Alfred Landé in 1921 \citep{Lande1921} fitted the anomalous Zeeman patterns
empirically with the formula
\begin{equation}\label{eq:lande_g}
  g_J = 1 + \frac{J(J+1) + S(S+1) - L(L+1)}{2J(J+1)},
\end{equation}
where $J$ is the total angular momentum quantum number and $L$, $S$ denote
orbital and ``spin'' contributions.  This formula reproduced the observed
line-splitting ratios remarkably well, but required $J$ to take \emph{half-integer}
values for many atomic states---impossible in the three-quantum-number scheme,
where all angular momenta are integers.

%======================================================================
\section{Stoner's 1924 Consolidation}
\label{sec:stoner}

The immediate empirical precursor to Pauli's paper was a 1924 article by
Edmund Clifton Stoner \citep{Stoner1924}.  Stoner used X-ray spectroscopic
level data (from Siegbahn and collaborators) to count the number of electrons
assigned to each sublevel of each shell.

Stoner's key observation: in an external magnetic field, all magnetic
sublevels ($m_l$ values) within a given $(n,l)$ subshell become energetically
distinct.  X-ray experiments showed that the number of electrons per subshell
is
\begin{equation}\label{eq:stoner}
  \text{capacity of subshell } (n,l) = 2(2l+1).
\end{equation}
The factor $2l+1$ counts the distinct $m_l$ values; each $m_l$ sublevel holds
\emph{two} electrons.  The mysterious factor of~2 per $m_l$ sublevel could
not be explained by any then-known quantum number.

Summing equation~\eqref{eq:stoner} over $l=0,1,\ldots,n-1$ gives
\begin{equation}\label{eq:shell_cap}
  \text{capacity of shell } n = \sum_{l=0}^{n-1} 2(2l+1)
  = 2\sum_{l=0}^{n-1}(2l+1) = 2n^2.
\end{equation}
This matches the observed noble-gas configurations exactly
(Table~\ref{tab:three_param}).  But Stoner did not explain \emph{why} each
$m_l$ sublevel holds two electrons.  Pauli read Stoner's paper and recognised
the answer: each sublevel is doubly occupied because there exists a fourth,
two-valued quantum number that takes two distinct values.  Occupancy is then
one electron per distinct 4-tuple, and the factor of~2 follows.

%======================================================================
\section{Pauli's Four-Parameter Model}
\label{sec:model}

\subsection{The Four Quantum Numbers}

Pauli's 1925 paper \citep{Pauli1925} introduces the following model.  Each
electron in an atom is completely characterised by four quantum numbers:

\begin{enumerate}
  \item \textbf{Principal quantum number} $n \in \{1,2,3,\ldots\}$:
    determines the shell and (for hydrogen) the energy $E_n \propto -1/n^2$.
    Range: any positive integer.

  \item \textbf{Azimuthal quantum number} $l \in \{0,1,\ldots,n-1\}$:
    determines the subshell and the orbital angular momentum magnitude.
    Range: $n$ possible values for each $n$.

  \item \textbf{Magnetic quantum number} $m_l \in \{-l,-l{+}1,\ldots,+l\}$:
    projection of orbital angular momentum on the field axis; determines the
    normal Zeeman sublevel.  Range: $2l+1$ possible values for each~$l$.

  \item \textbf{Fourth quantum number} $m_s \in \{+\tfrac{1}{2},
    -\tfrac{1}{2}\}$: a new, purely quantum-mechanical, two-valued
    property of the electron with no classical mechanical
    analogue.  Pauli called it a \emph{Zweideutigkeit}---two-valuedness.
    Range: exactly 2 values.
\end{enumerate}

\noindent Pauli could not give a physical interpretation for $m_s$ in 1925; he
was explicit that it is ``not classically describable.''  The physical
interpretation (intrinsic electron spin) was supplied a few months later by
Uhlenbeck and Goudsmit \citeyearpar{UhlenbeckGoudsmit1925}.

\subsection{The Exclusion Constraint}

With the four quantum numbers defined, the Exclusion Principle states:

\begin{center}
\fbox{\parbox{10cm}{%
\textbf{Pauli's Exclusion Principle (1925).}
In a multi-electron atom, no two electrons may simultaneously have identical
values of all four quantum numbers $(n, l, m_l, m_s)$.\vspace{-4pt}}}
\end{center}

\noindent Equivalently, each distinct 4-tuple $(n,l,m_l,m_s)$ is a quantum
\emph{state}, and every state may be occupied by at most one electron.
The occupation number of state $s$ satisfies
\begin{equation}\label{eq:occupation}
  n_s \in \{0,\,1\}.
\end{equation}
This binary constraint is the statistical mechanics analogue of a $\{0,1\}$ dummy
variable.  It is the defining property of \emph{fermions}---particles that obey
Fermi--Dirac statistics---in contrast to \emph{bosons}, for which occupation
numbers are unrestricted non-negative integers.

\subsection{Counting States: The Model's Predictions}

\paragraph{States per subshell.}
For fixed $(n,l)$, the number of distinct states is
\begin{equation}\label{eq:subshell_count}
  \underbrace{(2l+1)}_{\text{values of }m_l} \times
  \underbrace{2}_{\text{values of }m_s}
  = 2(2l+1).
\end{equation}
Under the Exclusion Principle, the maximum number of electrons in the
$(n,l)$ subshell equals the number of distinct states:
\begin{equation}\label{eq:subshell_cap}
  \text{capacity}(n,l) = 2(2l+1).
\end{equation}
Table~\ref{tab:subshell} lists the first four subshell types.

\begin{table}[ht]
\centering
\caption{Predicted subshell capacities from Pauli's four-parameter model,
equation~\eqref{eq:subshell_cap}.  Each entry agrees exactly with
observation.}
\label{tab:subshell}
\medskip
\begin{tabular}{ccccc}
\toprule
Subshell & $l$ & $m_l$ values & States $2(2l+1)$ & Observed electrons \\
\midrule
$s$ & 0 & $\{0\}$              &  2 &  2 \\
$p$ & 1 & $\{-1,0,+1\}$        &  6 &  6 \\
$d$ & 2 & $\{-2,-1,0,+1,+2\}$ & 10 & 10 \\
$f$ & 3 & $\{-3,-2,-1,0,+1,+2,+3\}$ & 14 & 14 \\
\bottomrule
\end{tabular}
\end{table}

\paragraph{States per shell.}
For fixed $n$, summing over subshells:
\begin{equation}\label{eq:shell_count}
  \text{capacity}(n) = \sum_{l=0}^{n-1} 2(2l+1)
  = 2\sum_{l=0}^{n-1}(2l+1) = 2n^2.
\end{equation}
\textit{Derivation:}\quad $\displaystyle\sum_{l=0}^{n-1}(2l+1)
  = 2\cdot\frac{(n-1)n}{2} + n = n^2$, so the shell capacity is $2n^2$.

Table~\ref{tab:shell} verifies these predictions against noble-gas
electron configurations.  The agreement is exact at every shell.

\begin{table}[ht]
\centering
\caption{Shell capacities predicted by Pauli's model vs.\ observed
electron configurations of noble gases.  ``Observed'' electron count is
the total for the filled inner shells, confirmed by X-ray spectroscopy and
the chemical inertness of the noble gases.}
\label{tab:shell}
\medskip
\begin{tabular}{ccccl}
\toprule
Shell & $n$ & Predicted $2n^2$ & Observed & Noble gas \\
\midrule
K & 1 &  2 &  2 & He ($Z=2$): $1s^2$ \\
L & 2 &  8 &  8 & Ne ($Z=10$): He $+\;2s^2 2p^6$ \\
M & 3 & 18 & 18 & Ar ($Z=18$): Ne $+\;3s^2 3p^6$ \\
N & 4 & 32 & 32 & Kr ($Z=36$): Ar $+\;4s^2 3d^{10} 4p^6$ \\
\bottomrule
\end{tabular}
\end{table}

%======================================================================
\section{The Empirical Dataset}
\label{sec:data}

The ``data'' that Pauli was required to organise fall into three classes,
each contributing a distinct type of \emph{empirical count} or
\emph{empirical multiplicity}.

\subsection{Class~1: Shell and Subshell Electron Capacities}

The electron capacities of each shell and subshell were known empirically
from two sources:
\begin{enumerate}
  \item \emph{Chemical valence and noble-gas stability.}  Bohr in 1922
    inferred from the periodic table that the noble-gas configurations
    correspond to closed shells with 2, 8, 18, 32 electrons.
  \item \emph{X-ray spectroscopy.}  Manne Siegbahn's group measured the
    binding energies and number of energy sublevels in the K, L, M, and N
    shells, identifying how many electrons reside at each distinct energy.
    Stoner \citeyearpar{Stoner1924} systematised these in his 1924 paper.
\end{enumerate}

The resulting empirical dataset is given in Table~\ref{tab:empirical_data}.
Each row is one ``observation.''

\begin{table}[ht]
\centering
\caption{Empirical electron capacities used by Pauli.  Shell and subshell
electron counts from chemical data (noble-gas configurations) and X-ray
spectroscopy (Siegbahn; systematised by \citealt{Stoner1924}).
These are the observables that Pauli's model must predict exactly.}
\label{tab:empirical_data}
\medskip
\begin{tabular}{ccccc}
\toprule
Shell & $n$ & Subshell & $l$ & Observed electrons \\
\midrule
K & 1 & $1s$ & 0 & 2 \\
L & 2 & $2s$ & 0 & 2 \\
  &   & $2p$ & 1 & 6 \\
M & 3 & $3s$ & 0 & 2 \\
  &   & $3p$ & 1 & 6 \\
  &   & $3d$ & 2 & 10 \\
N & 4 & $4s$ & 0 & 2 \\
  &   & $4p$ & 1 & 6 \\
  &   & $4d$ & 2 & 10 \\
  &   & $4f$ & 3 & 14 \\
\bottomrule
\end{tabular}
\end{table}

\subsection{Class~2: Spectral Term Multiplicities}

The \emph{multiplicity} of a spectral term ${}^{2S+1}L_J$ counts the number
of distinct quantum states with a given $L$, $S$, and $J$: specifically, it
equals $2J+1$ (the number of $m_J$ projections).  The observed term
multiplicities for the first few elements are:

\begin{itemize}
  \item \emph{Alkali metals} (Li, Na, K, \ldots): doublet spectra
    (${}^2S_{1/2}$, ${}^2P_{1/2}$, ${}^2P_{3/2}$, \ldots).
    The doublet character---two lines everywhere the classical single-electron
    theory predicted one---was a direct signal of $S = \tfrac{1}{2}$.
  \item \emph{Alkaline earths} (Be, Mg, Ca, \ldots): singlet and triplet
    series (${}^1S_0$, ${}^3S_1$, ${}^1P_1$, ${}^3P_{0,1,2}$, \ldots).
  \item \emph{Carbon, nitrogen, oxygen}: quartet terms among excited
    configurations.
\end{itemize}

\noindent These term multiplicities are predicted exactly by the
four-quantum-number scheme: for a closed subshell the total spin is zero
(singlet); for an outer shell with one electron the total spin is $S=\tfrac{1}{2}$
(doublet); for two valence electrons in different orbitals, $S \in \{0,1\}$
(singlet or triplet).

\subsection{Class~3: Anomalous Zeeman Component Counts}

In a weak external magnetic field, each spectral term ${}^{2S+1}L_J$ splits
into $2J+1$ sublevels.  A transition ${}^{2S+1}L_J \to {}^{2S'+1}L'_{J'}$
then produces a number of distinct emission lines equal to the number of
allowed ($|\Delta m_J| \leq 1$, $\Delta m_J \in \{-1,0,+1\}$) transitions
between those sublevels.  Table~\ref{tab:zeeman_data} lists key observed
anomalous Zeeman patterns.

\begin{table}[ht]
\centering
\caption{Observed anomalous Zeeman component counts for selected sodium
spectral lines.  The upper and lower $J$ values and the corresponding
$m_J$-sublevel counts are listed.  ``Predicted'' uses the four-quantum-number
model with the exclusion principle and is identical to the observed count in
every case.  Classical Lorentz theory always predicts 3 components.}
\label{tab:zeeman_data}
\medskip
\begin{tabular}{lccccc}
\toprule
Transition &
Line (nm) &
Upper $J$ &
Lower $J$ &
Observed lines &
Classical \\
\midrule
$^2P_{3/2}\to{}^2S_{1/2}$ & 589.0 (D$_2$) & 3/2 & 1/2 & 6 & 3 \\
$^2P_{1/2}\to{}^2S_{1/2}$ & 589.6 (D$_1$) & 1/2 & 1/2 & 4 & 3 \\
\bottomrule
\end{tabular}
\end{table}

\noindent \emph{Calculation for the D$_2$ line.}
The upper state $^2P_{3/2}$ has $J=3/2$, giving $m_J \in
\{-3/2,-1/2,+1/2,+3/2\}$ (4~sublevels).  The lower state $^2S_{1/2}$
has $J=1/2$, giving $m_J\in\{-1/2,+1/2\}$ (2~sublevels).  The allowed
transitions with $\Delta m_J \in \{-1,0,+1\}$ number exactly~6, as observed.
For the D$_1$ line, both levels have $J=1/2$; the 4 allowed transitions
are $(\Delta m_J = 0, \pm1)$ from 2 upper to 2 lower sublevels
(one $\Delta m_J = 0$ is forbidden for $J=1/2\to1/2$, $\Delta J=0$
by the selection rule, leaving~4).~\footnote{Precisely: for $\Delta J = 0$
the transition $m_J=0 \to m_J=0$ is allowed only if $J \neq 0$.  For
the D$_1$ line the two $\Delta m_J=0$ transitions both occur ($+1/2\to+1/2$
and $-1/2\to-1/2$), giving 2 longitudinal lines, combined with 1 transverse
$\sigma^+$ and 1 $\sigma^-$, totalling~4.}

These counts---6 and 4 instead of the classical 3---are \emph{impossible}
within the three-quantum-number scheme, which yields only integer $J$ and
hence cannot produce the half-integer $J=3/2$ or $J=1/2$ occurring in
the sodium D~doublet.

\subsection{Landé's g-Factor Data}

Landé's 1921 empirical formula~\eqref{eq:lande_g} summarises a further
dataset: the \emph{spacing ratios} of Zeeman components.  The splitting
of a level in a magnetic field $B$ is $\Delta E = g_J \mu_B m_J B$, where
$\mu_B = e\hbar/(2m_e)$ is the Bohr magneton.  Observed $g_J$ values
for sodium are:
\begin{align*}
  g_J({}^2S_{1/2}) &= 2.0, \\
  g_J({}^2P_{1/2}) &= 2/3 \approx 0.667, \\
  g_J({}^2P_{3/2}) &= 4/3 \approx 1.333.
\end{align*}
These non-integer values---inexplicable if $g_{\mathrm{spin}} = 1$ as classical
orbital mechanics would have it---require $g_s = 2$ for the spin contribution.
Formula~\eqref{eq:lande_g} with $g_s = 2$ reproduces all observed $g_J$
values exactly.  The anomalous value $g_s = 2$ is identified by the Zeeman
spacing data.

%======================================================================
\section{Goodness of Fit: Zero Residuals}
\label{sec:fit}

Table~\ref{tab:fit} compares Pauli's predicted electron capacities with
every observed capacity from Table~\ref{tab:empirical_data}.

\begin{table}[ht]
\centering
\caption{Predicted vs.\ observed subshell electron capacities.  Pauli's
model predicts $2(2l+1)$ electrons per subshell $(n,l)$.  The residual
is zero in every case; the model fits all ten observations exactly with
no adjustable parameters beyond the ranges of the four quantum numbers.}
\label{tab:fit}
\medskip
\begin{tabular}{cccrrc}
\toprule
Subshell & $l$ & Formula $2(2l+1)$ &
Predicted & Observed & Residual \\
\midrule
$1s$ & 0 &  $2(1)$  &  2 &  2 & $0$ \\
$2s$ & 0 &  $2(1)$  &  2 &  2 & $0$ \\
$2p$ & 1 &  $2(3)$  &  6 &  6 & $0$ \\
$3s$ & 0 &  $2(1)$  &  2 &  2 & $0$ \\
$3p$ & 1 &  $2(3)$  &  6 &  6 & $0$ \\
$3d$ & 2 &  $2(5)$  & 10 & 10 & $0$ \\
$4s$ & 0 &  $2(1)$  &  2 &  2 & $0$ \\
$4p$ & 1 &  $2(3)$  &  6 &  6 & $0$ \\
$4d$ & 2 &  $2(5)$  & 10 & 10 & $0$ \\
$4f$ & 3 &  $2(7)$  & 14 & 14 & $0$ \\
\midrule
\multicolumn{4}{l}{RSS} & & $0$ \\
\bottomrule
\end{tabular}
\end{table}

\noindent The residual sum of squares is identically zero.  A formal
$\chi^2$ goodness-of-fit test statistic
\[
  \chi^2 = \sum_{\text{subshells}}
  \frac{(\text{Observed} - \text{Predicted})^2}{\text{Predicted}}
  = 0.
\]
Pauli's model does not fit the data approximately; it generates the
observed counts \emph{deductively}.  The fitting exercise has zero
degrees of freedom in the residual because the model produces
the observed counts as exact consequences of counting arguments,
not as the output of a curve-fit with estimated free parameters.

\medskip
This situation precisely parallels Einstein's photoelectric equation: the
five-point OLS regression for sodium gave $R^2 \approx 0.99999$ and
residuals below 2\,mV, but crucially the \emph{functional form}
$V_0 = \alpha + \beta\nu$ and the \emph{universality of the slope} $\beta = h/e$
were \emph{predicted} by theory, not estimated from data.  For Pauli, the
functional form of the capacity formula $2(2l+1)$ is similarly predicted by
theory; the counting data serve as a confirmation test, not as inputs to a
parameter-estimation algorithm.

%======================================================================
\section{Parameter Identification}
\label{sec:identification}

In a classical regression, ``identification'' asks which variation in the data
pins down which parameter.  For Pauli's lattice of allowed states, the
analogous question is: which empirical observation establishes each of the four
quantum numbers and the exclusion constraint?

\begin{center}
\begin{tabular}{p{3.8cm}p{5.5cm}p{4.5cm}}
\toprule
Parameter & What it determines & Identifying observation \\
\midrule
$n \in \{1,2,3,\ldots\}$ & Shell index; number of subshells
  $= n$ & Chemical periodicity; discrete X-ray absorption edges \\
$l \in \{0,\ldots,n-1\}$ & Subshell type; $2l+1$ magnetic sublevels
  & Fine structure of X-ray levels; Sommerfeld's elliptical orbits \\
$m_l \in \{-l,\ldots,+l\}$ & Normal Zeeman splitting into $2l+1$ components
  & Normal Zeeman effect (integer-valued splitting,
    $\Delta\nu \propto m_l$) \\
$m_s \in \{\pm\tfrac{1}{2}\}$ & Factor-of-2 per $m_l$ sublevel;
  doublet structure in alkali spectra
  & Stoner's factor-of-2 per X-ray sublevel; alkali-metal
    doublets; Landé $g_s = 2$ \\
Exclusion ($n_s \leq 1$) & Total shell capacity $= 2n^2$;
  chemical inertness of noble gases
  & Noble-gas stability; absence of transitions out of
    closed shells \\
\bottomrule
\end{tabular}
\end{center}

\paragraph{Remark on the exclusion constraint.}
Without the constraint $n_s \leq 1$, multiple electrons could pile into the
ground state~$1s$ (as bosons do).  Every atom would have all its electrons in
the $n=1$, $l=0$, $m_l=0$, $m_s=+\tfrac{1}{2}$ state, and the periodic
table would not exist.  The chemical diversity of matter is a consequence
solely of this constraint.  In statistical language, the constraint is
identified by the existence of the periodic table itself: the chemical fact
that atoms do not all collapse to a single ground-state configuration is
direct evidence that electrons must obey an exclusion rule.

\paragraph{Remark on $m_s$.}
The identification of $m_s$ is notably \emph{indirect}: the empirical evidence
for a fourth quantum number comes from its \emph{consequences} (the factor of~2
in subshell capacities; the doublet line structure) rather than from a direct
measurement of the quantum number itself as a physical observable.  Pauli
acknowledged this explicitly in his Nobel lecture: ``I was unable to give a
logical reason for the exclusion principle or to derive it from more general
assumptions'' \citep{Pauli1946Nobel}.  The physical object
corresponding to $m_s$ was identified only after Uhlenbeck and Goudsmit proposed
electron spin.

%======================================================================
\section{A Maximum Likelihood Framing}
\label{sec:mle}

The photoelectric regression has a natural MLE formulation under Gaussian
errors.  Pauli's counting problem is discrete, and its counterpart is a
multinomial likelihood.  We set this up here to make the statistical parallel
explicit.

\subsection{The Discrete Probability Model}

Consider the experiment of drawing an electron uniformly at random from among
all quantum states available in shell~$n$, and recording its $(l, m_l, m_s)$
values.  Under the Exclusion Principle, each available state is equally
accessible (the Pauli principle says nothing about energy ordering; for the
purpose of counting, all states are equivalent).  The relevant statistical
object is therefore the \emph{marginal probability} that the electron is in
subshell~$l$:
\begin{equation}\label{eq:marginal}
  p(l \mid n) = \frac{2(2l+1)}{2n^2}
              = \frac{2l+1}{n^2},
  \quad l=0,1,\ldots,n-1.
\end{equation}
This is a discrete distribution on $\{0,1,\ldots,n-1\}$ with probabilities
proportional to the subshell degeneracy $2l+1$.

\subsection{The Multinomial Likelihood}

Suppose we observe $N_l$ electrons in subshell~$l$ for $l=0,1,\ldots,n-1$,
with $\sum_l N_l = N$ (the total number of electrons in shell~$n$ at
maximum fill).  The multinomial likelihood is
\begin{equation}\label{eq:multinomial}
  \mathcal{L}(\{p_l\}) \propto \prod_{l=0}^{n-1} p_l^{N_l}.
\end{equation}
Pauli's model predicts the \emph{unique} probability vector
$p_l^* = (2l+1)/n^2$.  Any competing model that, for example, assigns equal
probability to all subshells ($p_l = 1/n$) or uses only three quantum numbers
($p_l = 1/(n^2)$ per $m_l$ if no $m_s$) would predict different capacities.

Under the true multinomial, the log-likelihood ratio comparing Pauli's model
to the null model of equal subshell occupancy ($p_l^{\text{null}} = 1/n$ for all $l$)
is
\begin{equation}\label{eq:llr}
  \Lambda = 2\sum_{l=0}^{n-1} N_l \ln\!\frac{p_l^*}{p_l^{\text{null}}}
          = 2N \sum_{l=0}^{n-1} \frac{2l+1}{n^2}\ln\!\frac{n(2l+1)}{n^2},
\end{equation}
which is strictly positive (Pauli's model fits better) for all $n\geq 2$.
The $\chi^2$ statistic comparing the predicted capacities from
Table~\ref{tab:fit} with the null prediction of $2n^2/n$ electrons per
subshell is:

\begin{center}
\begin{tabular}{ccccr}
\toprule
Shell $n$ & Subshell $l$ & Predicted ($2(2l+1)$) &
  Null ($2n^2/n = 2n$) & $\chi^2$ contribution \\
\midrule
$n=2$ & $l=0$ ($s$) &  2 & 4 & $4/4=1.00$ \\
$n=2$ & $l=1$ ($p$) &  6 & 4 & $4/4=1.00$ \\
\midrule
\multicolumn{4}{r}{Shell $n=2$ total $\chi^2$} & $2.00$ \\
\midrule
$n=3$ & $l=0$ ($s$) &  2 & 6 & $16/6=2.67$ \\
$n=3$ & $l=1$ ($p$) &  6 & 6 & $0/6=0.00$ \\
$n=3$ & $l=2$ ($d$) & 10 & 6 & $16/6=2.67$ \\
\midrule
\multicolumn{4}{r}{Shell $n=3$ total $\chi^2$} & $5.33$ \\
\bottomrule
\end{tabular}
\end{center}

\noindent The null model of equal subshell occupancy is clearly rejected by
the observed capacity data at any conventional significance level.  Pauli's
model, by contrast, fits perfectly.

\subsection{Parsimony: The Model Has No Free Parameters}

The multinomial likelihood for the capacity data has, in principle, $n-1$
free probability parameters for each shell~$n$.  Pauli's model reduces these
to \emph{zero} free parameters: the probabilities $p_l^* = (2l+1)/n^2$ are
completely determined by the structure of the four quantum numbers and the
counting rule.  There is nothing to estimate.

In the language of information criteria, if a competitor model uses $k$ free
parameters and Pauli's model uses~0, the Akaike Information Criterion (AIC)
favours Pauli's model by $2k$ units \emph{plus} the improvement in
log-likelihood from the exact fit.  Since the exact fit achieves the maximum
possible likelihood (the observed data are the maximum-likelihood estimate
of their own probabilities), no model can outpredict Pauli's on data
that were used to derive the principle --- but Pauli's model generalises
analytically to any shell, including $n=5$ (xenon) and $n=7$ (actinides),
whereas a purely empirical fit to shells $n=1,\ldots,4$ would have no
basis for predicting shell~5.

%======================================================================
\section{Subsequent Confirmations}
\label{sec:subsequent}

\subsection{Uhlenbeck and Goudsmit: Physical Interpretation of $m_s$}

Pauli proposed $m_s$ as a purely formal device in January 1925.  In
September 1925, George Uhlenbeck and Samuel Goudsmit
\citeyearpar{UhlenbeckGoudsmit1925} proposed that the electron possesses an
\emph{intrinsic angular momentum} (spin) of magnitude $\hbar/2$, with
projection $m_s = \pm\tfrac{1}{2}$ along any axis.  The associated magnetic
moment has an anomalous gyromagnetic ratio $g_s \approx 2$---twice the
classical orbital value.  The Uhlenbeck--Goudsmit proposal:
\begin{itemize}
  \item Gave a physical object---intrinsic spin---to Pauli's abstract
    two-valuedness.
  \item Explained the doublet structure of alkali spectra as arising from the
    spin--orbit coupling $\mathbf{L}\cdot\mathbf{S}$ between orbital and spin
    angular momenta.
  \item Reproduced the Landé g-factor formula~\eqref{eq:lande_g} with the
    substitution $g_s = 2$, explaining why the formula contains the factor~2
    in the numerator.
\end{itemize}

\subsection{Fermi--Dirac Statistics}

In 1926 Enrico Fermi \citeyearpar{Fermi1926} and Paul Dirac
\citeyearpar{Dirac1926} independently derived the statistical mechanics of
a gas of identical particles obeying the Exclusion Principle.  For a system
in thermal equilibrium at temperature $T$ with chemical potential $\mu$, the
mean occupation number of a state with energy $\varepsilon$ is
\begin{equation}\label{eq:fermi_dirac}
  \bar{n}(\varepsilon) = \frac{1}{e^{(\varepsilon-\mu)/k_BT} + 1}.
\end{equation}
The $+1$ in the denominator (vs.\ the $-1$ of Bose--Einstein statistics)
is the signature of the Exclusion Principle: the constraint $n_s\leq 1$
prevents accumulation, and $0 < \bar{n} < 1$ for all $\varepsilon$ and $T>0$.
At $T=0$, equation~\eqref{eq:fermi_dirac} becomes a step function: all states
below the Fermi energy $E_F = \mu(T=0)$ are completely filled
($\bar{n}=1$) and all above are empty ($\bar{n}=0$)---a direct restatement
of the Exclusion Principle as a thermodynamic limiting case.

Fermi--Dirac statistics had immediate astrophysical consequences.
\citet{Fowler1926} applied them to explain the stability of white dwarf
stars: the electron degeneracy pressure arising from the Exclusion Principle
resists gravitational collapse even at zero temperature.  Sommerfeld
\citeyearpar{Sommerfeld1927} resolved the long-standing puzzle of why
conduction electrons contribute far less to metallic heat capacity than
classical equipartition predicts: at room temperature, $k_BT \ll E_F$,
and the Pauli exclusion principle suppresses thermal excitations for all
but the small fraction of electrons within $k_BT$ of the Fermi surface.

\subsection{Dirac's Relativistic Derivation}

Paul Dirac derived the relativistic quantum equation for the electron in 1928
\citep{Dirac1928}.  The Dirac equation requires the wave function to be a
four-component spinor, automatically producing:
\begin{itemize}
  \item Intrinsic spin-$\tfrac{1}{2}$: the spin degree of freedom is not
    added by hand but emerges from the algebra of the equation.
  \item $g_s = 2$ exactly: the anomalous gyromagnetic ratio of Uhlenbeck and
    Goudsmit is a direct consequence of the Dirac equation; quantum
    electrodynamic corrections produce the small deviation
    $g_s = 2.00232\ldots$
\end{itemize}
Dirac's derivation retroactively justified the Uhlenbeck--Goudsmit proposal
and gave Pauli's fourth quantum number its derivation from first principles.

\subsection{The Spin--Statistics Theorem}

Pauli proved in 1940 \citep{Pauli1940} that in any relativistic quantum field
theory satisfying locality, causality, and positivity of energy, particles with
half-integer spin \emph{must} obey the Exclusion Principle (Fermi--Dirac
statistics), while particles with integer spin must obey Bose--Einstein
statistics.  The Exclusion Principle is thus not an independent postulate in quantum
field theory: it is a theorem, a consequence of combining quantum mechanics
with special relativity in a consistent way.  This is the deepest confirmation
of Pauli's 1925 heuristic.

Pauli received the Nobel Prize in Physics in 1945 specifically ``for the
discovery of the Exclusion Principle, also called the Pauli Principle.''

%======================================================================
\section{Pauli as Machine Learner}
\label{sec:ml}

The vocabulary of modern machine learning---supervised learning, hypothesis
class, inductive bias, feature engineering, model selection---maps onto
Pauli's scientific reasoning in almost the same way it maps onto Einstein's
photoelectric reasoning.

\subsection{The Learning Problem}

The ``dataset'' assembled by Pauli is a table of discrete counts
(Table~\ref{tab:empirical_data}) and multiplicities
(Tables~\ref{tab:zeeman_data},~\ref{tab:fit}).  The task is to find a
model that predicts all entries in these tables from a minimal set of
structural rules.  In modern ML parlance, this is a multivariate
\emph{classification} or \emph{generative model} problem over a discrete
state space.

\subsection{The Hypothesis Class and Inductive Bias}

Einstein's hypothesis class was the set of affine functions
$f_{\alpha,\beta}:\nu\mapsto\alpha+\beta\nu$, derived from the light-quantum
postulate.  Pauli's hypothesis class is the set of binary occupation functions
on four-integer-labeled states, derived from two postulates:
\begin{enumerate}
  \item \emph{Four quantum numbers}: the state space is the Cartesian product
    $\{1,2,\ldots\}\times\{0,\ldots,n-1\}\times\{-l,\ldots,+l\}\times
    \{+\tfrac{1}{2},-\tfrac{1}{2}\}$.
  \item \emph{Exclusion}: occupation numbers are in $\{0,1\}$.
\end{enumerate}
Together these postulates restrict the capacity of each subshell to exactly
$2(2l+1)$---a prediction with zero adjustable parameters beyond the
specification of the state space.

Prior to Pauli, the effective hypothesis class was a much larger family: models
with arbitrary occupation rules, non-integer quantum numbers, or additional
continuous parameters.  Pauli's postulates function as an extreme inductive
bias that collapses all combinatorial freedom into a single prediction.  In
learning-theoretic terms, the VC dimension of Pauli's ``model class'' (the
lattice of binary-occupied states subject to the exclusion constraint) is
effectively~1: there is only one model in the class (given the four quantum
numbers and the exclusion rule, the capacity function is uniquely determined),
so any single observation that confirms the predicted capacity suffices to
``learn'' the model.

\subsection{Feature Selection}

Pauli's four quantum numbers are four \emph{features} of each electron's
state.  By 1924, three of these features were known from Sommerfeld's work;
the problem was that the known feature set $(n, l, m_l)$ was incomplete.
Pauli's scientific contribution was to recognise that a fourth binary feature
is needed---not because a large dataset demanded it, but because the observed
capacity data were off by the precise factor of~2 that a binary feature adds.

This is a form of \emph{feature engineering}: diagnosing that the existing
feature set is insufficient, inferring the minimum additional feature---a
two-valued one---needed to close the discrepancy, and verifying that the
augmented feature set fits all observations exactly.

\subsection{What a Naive Data-Driven Learner Would Miss}

A purely data-driven learning system confronting Table~\ref{tab:empirical_data}
would observe ten (subshell, count) pairs.  A regression tree or neural
network could be trained to predict these counts from $(n,l)$.  What it could
not do without theoretical structure:

\begin{enumerate}
  \item It would not know that the correct model has the form $2(2l+1)$---it
    would fit a table of numbers and perhaps find the pattern numerically.
  \item It would not know that the explanation of the pattern is a fourth quantum
    number with exactly two values.  The data are consistent with many mathematical
    formulas; the fourth quantum number is the physically and theoretically
    minimal explanation.
  \item It would not generalise to the anomalous Zeeman effect predictions
    without the physical model as a common underlying theory.  The connection
    between shell capacity data and Zeeman splitting data requires the quantum
    number structure as the shared latent variable.
  \item It could not predict the Fermi distribution~\eqref{eq:fermi_dirac} or
    white dwarf stability---applications that rely on the Exclusion Principle
    in a continuum of energies, far outside the discrete atomic-level dataset.
\end{enumerate}

Pauli's irreplaceable contribution was the theoretical unification: a single
two-sentence postulate that simultaneously explains the capacity data, the
spectral multiplicities, and the Zeeman patterns, and that generalises to
all subsequent applications of the Exclusion Principle in condensed matter
physics, nuclear physics, and quantum field theory.

\subsection{Transfer Learning}

Einstein's slope $h/e$ transferred exactly from sodium to lithium---a
zero-shot prediction confirmed to 0.2\%.  Pauli's shell-capacity formula
$2n^2$ transfers from atoms to molecules (Hückel's rule for aromatic
stability), to nuclei (nuclear shell model), to solid-state crystals
(Fermi surfaces), and to neutron stars (neutron degeneracy pressure).
Each application is zero-shot: the Exclusion Principle was not re-estimated
but transferred with no adjustment, constrained only by the energy levels
applicable in the new domain.

%======================================================================
\section{Comparison with the Photoelectric Episode}
\label{sec:comparison}

Table~\ref{tab:comparison} compares the two episodes side by side.

\begin{table}[ht]
\centering
\caption{Structural comparison of Einstein's photoelectric derivation (1905--1916)
and Pauli's exclusion principle derivation (1924--1925) from the perspective
of statistical inference and machine learning.}
\label{tab:comparison}
\medskip
\begin{tabular}{p{3.5cm}p{5cm}p{5cm}}
\toprule
Dimension & Einstein (photoelectric) & Pauli (exclusion principle) \\
\midrule
Empirical anomaly & Stopping potential
  independent of intensity; depends on $\nu$ &
  Shell capacity $2n^2$ not $n^2$; anomalous
  Zeeman multiplicities \\
Key precursor data & Lenard (1902):
  $V_0$ vs.\ $\nu$ qualitative &
  Stoner (1924): $2$ electrons per $m_l$
  sublevel \\
Theoretical postulate & Light quantum:
  $E = h\nu$ &
  Fourth quantum number $m_s$; exclusion
  constraint $n_s \leq 1$ \\
Hypothesis class & Affine functions:
  $V_0 = \alpha + \beta\nu$ &
  Binary occupation on four-dimensional
  integer lattice \\
Free parameters & 2 (slope $\beta = h/e$;
  intercept $\alpha = -W/e + v_c$) &
  0 (all capacities determined by counting) \\
Prediction & $V_0 = (h/e)\nu - W/e$;
  universal slope across metals &
  Capacity$(n,l) = 2(2l+1)$;
  universal across all atoms and shells \\
Fit statistic & OLS: $R^2 = 0.99999$,
  $\hat\sigma = 1.4$\,mV (Millikan 1916) &
  $\chi^2 = 0$: exact fit to all 10
  subshell observations \\
Transfer & Na and Li slopes agree to 0.2\% &
  Formula transfers to molecules, nuclei,
  stars \\
Physical interpretation delayed & Light
  quantum debated until Compton (1923) &
  $m_s$ uninterpreted until Uhlenbeck--Goudsmit
  (1925) \\
Deductive grounding achieved & QED:
  photoelectric equation exact &
  Dirac (1928) and Pauli (1940) spin--statistics
  theorem \\
Nobel Prize & 1921 (Einstein, for
  photoelectric equation) &
  1945 (Pauli, for exclusion principle) \\
\bottomrule
\end{tabular}
\end{table}

A key structural difference is the \emph{type of data}.  Einstein's equation
is a linear regression on continuous data $(\nu_i, V_{0i})$; the free
parameters $\beta = h/e$ and $\alpha = -W/e + v_c$ are estimated by
OLS/MLE with degrees of freedom from three or more observations.  Pauli's
model involves a discrete counting argument; the ``free parameters'' are the
integer ranges of $n$, $l$, $m_l$, and the binary range of $m_s$, and they
are identified by qualitative and semi-quantitative spectroscopic data rather
than by a quantitative regression.

A second structural difference is the \emph{role of error}.  Einstein's
regression has a genuine residual $\varepsilon$ from measurement noise;
Millikan's precision was finite and the standard error of $h/e$ was
estimated at $0.1\%$.  Pauli's counting formula predicts integers; there is
no room for statistical noise---the capacity of the $2p$ subshell is exactly~6,
not $6.1 \pm 0.3$.  The ``goodness of fit'' is either perfect or it is not.

In both cases, theory supplied the hypothesis class before quantitative
data were available: Einstein wrote down $V_0 = (h/e)\nu - W/e$ in 1905
before clean data existed; Pauli wrote down the exclusion rule and the
fourth quantum number in January 1925, relying principally on the qualitative
patterns in Stoner's (1924) table.  In both cases, the theoretical prior
reduced the effective number of free parameters to near zero, enabling
confirmation with a minimal dataset.

%======================================================================
\section{Pauli's Epistemic Situation}

Pauli's relationship to his own principle shares certain features with
Millikan's relationship to Einstein's equation.  As Millikan confirmed
Einstein's equation while for years publicly rejecting its physical
interpretation, Pauli proposed the exclusion principle and the fourth quantum
number while explicitly refusing to interpret $m_s$ physically:

\begin{quote}
  ``Already in my original paper I stressed the circumstance that I was
  unable to give a logical reason for the exclusion principle or to derive
  it from more general assumptions\ldots\ To me, the suggestion of Stoner
  appeared to be an important step forward, since for the first time the
  number~2, characteristic of the closed shells, had to be attributed to the
  electron itself---to its `two-valuedness.'\,'' \citep{Pauli1946Nobel}
\end{quote}

The lesson is the same as in the photoelectric case: a sufficiently parsimonious
formal model can unify a body of empirical observations before the physical
mechanism underlying the model is understood.  Pauli in 1925 was, in modern
parlance, fitting a discrete probability model over quantum state space; the
physical object that corresponds to $m_s$---intrinsic spin---was supplied by
Uhlenbeck and Goudsmit a few months later, and its relativistic necessity was
established by Dirac three years after that.

When Uhlenbeck and Goudsmit proposed electron spin, Pauli was initially
sceptical, pointing out that a classically spinning electron the size of a
point particle would need a surface velocity exceeding the speed of light.
This reservation was eventually resolved by Dirac's equation, in which spin
is a purely quantum-mechanical property with no classical analogue---exactly
the ``classically non-describable two-valuedness'' Pauli had insisted upon
from the start.

%======================================================================
\section{Conclusions}
\label{sec:conclusions}

Pauli's 1925 determination of the quantum-mechanical state space for atomic
electrons is a canonical exercise in scientific model specification.  Starting
from a handful of empirical regularities---Stoner's factor-of-2 per sublevel,
the doublet structure of alkali spectra, the anomalous Zeeman multiplicities,
and the noble-gas closure counts---Pauli inferred the existence of a fourth,
two-valued quantum number and a binary occupation constraint.  His model has
the structure of a four-parameter discrete probability distribution over a
lattice, with the unusual property that the parameters are integers (the
ranges of $n$, $l$, $m_l$, $m_s$) and the predictions are exact counts,
leaving zero tunable parameters and zero residual degrees of freedom.

Three features of the episode are especially instructive from a statistical
standpoint.

\medskip
\emph{First, extreme parsimony.}  The principle is a single sentence
($n_s \leq 1$), yet it organises the structure of the entire periodic
table, explains all atomic spectral multiplicities, predicts the anomalous
Zeeman patterns, and implies Fermi--Dirac statistics.  The
Akaike-information-criterion premium for zero free parameters---relative
to any competitor that uses tunable constants---is overwhelmingly in Pauli's
favour.

\medskip
\emph{Second, theory-first identification.}  Like Einstein's photoelectric
paper, Pauli's principle was proposed before quantitative experimental
confirmation was complete.  The fourth quantum number was inferred from
the inability of three quantum numbers to explain a well-known discrepancy
(factor of~2), not from a new precision measurement.  The measurement
confirmation---via Zeeman pattern counts, X-ray level counts, and chemical
periodicities---served to verify a pre-existing theoretical prediction.

\medskip
\emph{Third, uninterpreted success.}  Pauli in 1925 did not know what $m_s$
was.  In this he resembles Planck in 1900, who introduced the energy quantum
$h\nu$ to fit the blackbody spectrum without knowing what a ``quantum''
physically was, and Einstein in 1905, who proposed the light quantum while
acknowledging that it conflicted with the wave theory of light.  Each of these
model specifications---Planck's, Einstein's, Pauli's---introduced a formal
device that organised empirical data parsimoniously, long before a
self-consistent physical interpretation was available.  The pattern suggests
that in fundamental physics, as in econometrics, the ability to specify a
parsimonious model is often the critical intellectual step, and the ability to
interpret its parameters often follows much later.

\bigskip
\bibliographystyle{plainnat}
\bibliography{pauli_exclusion_ML}

\end{document}
